\honest{Scarcity}{Eliezer Yudkowsky}{2008}

What follows comes mostly from Robert Cialdini's \textsc{Influence}. I own three copies, two of them for lending.

Scarcity, in social psychology, is when things become more desirable as they appear less obtainable. Two-year-old boys go straight for a toy behind a tall Plexiglas wall, and show no preference when the wall is low. \nb{Cialdini's endnote adds that the girls in the study did not respond this way.} After Dade County banned phosphate detergents, its residents drove to other counties to stock up, and rated the detergents better than people in Tampa did.

Information works the same way. Students who learned that a speech against coed dorms had been banned grew more opposed to coed dorms. I cite it as ``Probably in Ashmore et. al. 1971.'' \nb{A 1971 paper by Worchel and Arnold describes Ashmore's study as one about police on college campuses. Another study Cialdini cites found that when an attractive expert censor banned a speech the audience disagreed with, interest in it fell.} Jurors told that a driver was insured awarded \$4,000 more; told also to disregard the insurance, they awarded \$13,000 more. \nb{A 2006 meta-analysis found that instructions to disregard evidence fail to remove its effect; its abstract does not say they generally increase it.} Supermarket buyers told beef would be scarce ordered twice as much, and six times as much when told the news was exclusive. \nb{The comparison group heard the usual sales pitch, not that beef was plentiful. The study is an unpublished dissertation by Cialdini's student, who owned the company.} I add that the information was in fact correct. \nb{Cialdini confirms it: there was a real shortage. So these buyers ordered more on true news of a coming shortage, and the study does not show that they erred.}

The usual explanation is ``psychological reactance'': tell people they can't do something and they try harder. We resist anyone who restricts our freedom, and we seize options before they vanish. \nb{Cialdini names a first cause before reactance: scarce things are usually better, so availability is a quick guide to quality, and ``by following it, we are usually and efficiently right.'' The post leaves this out.} Seizing vanishing options may have suited hunter-gatherers, I say, but with money it is costly. An appliance salesman tells interested customers the last one was sold twenty minutes ago, then offers to check the back room if they will commit to buy.

Cialdini's sign of ``this malfunction,'' I say, is dreaming of possessing a thing rather than using it. \nb{Cialdini says that if you want a rare thing for the benefits of owning it, scarcity ``will give us a good indication'' of its value. The warning is for things you want to use.}

The deepest problem is that what you desire because it is unattainable stops being unattainable once you have it. ``If we cannot take joy in the merely available, our lives will always be frustrated.''
